% DiVA -- bản thảo theo template ACL (acl-org/acl-style-files).
% Biên dịch bằng pdfLaTeX: latexmk -pdf diva.tex
% (cần gói vntex cho bảng mã T5 tiếng Việt; Overleaf có sẵn).
\documentclass[11pt]{article}

% "review" = bản nộp ẩn danh; "preprint" = bản có tên tác giả và số trang; "final" = bản camera-ready.
\usepackage[review]{acl}

\usepackage{times}
\usepackage{latexsym}
\usepackage[T5]{fontenc}   % tiếng Việt (theo hướng dẫn của template ACL)
\usepackage[utf8]{inputenc}
\usepackage{microtype}
\usepackage{inconsolata}
\usepackage{graphicx}

\usepackage{amsmath}
\usepackage{amssymb}
\usepackage{booktabs}
\usepackage{array}
\newcolumntype{L}[1]{>{\raggedright\arraybackslash}p{#1}}
\usepackage{multirow}
\usepackage{enumitem}
\usepackage{tikz}
\usetikzlibrary{arrows.meta,positioning,fit,backgrounds,shapes.geometric}

\setlist{nosep,leftmargin=*}
\renewcommand{\abstractname}{Tóm tắt}
\renewcommand{\figurename}{Hình}
\renewcommand{\tablename}{Bảng}
\newcommand{\pchot}{p_{\text{chốt}}}
\newcommand{\diva}{\textsc{DiVA}}
\newcommand{\code}[1]{{\fontencoding{T1}\selectfont\texttt{#1}}}

\title{DiVA: Sinh biên bản họp tiếng Việt có căn cứ\\
bằng cắt chủ đề hội thoại và đa tác tử kiểm chứng}

\author{Ẩn danh \\
  \texttt{anonymous@example.com}}

\begin{document}
\maketitle

\begin{abstract}
Ở cơ quan nhà nước, doanh nghiệp hay nhóm dự án, cuộc họp là nơi ra quyết định và phân công công việc.
Biên bản họp vì thế là căn cứ về \emph{trách nhiệm}: ai làm gì, phối hợp với ai, hạn khi nào, đã thống
nhất điều gì, và việc giao lần trước đã tiến triển ra sao. Một hệ thống tự động soạn biên bản phải đúng
về những điểm này, không chỉ trôi chảy về câu chữ. Chúng tôi phát biểu bài toán \emph{sinh biên bản họp
có căn cứ} từ bản ghi giọng nói tiếng Việt, gồm bốn phần: diễn biến, giao việc, kết luận và tiến độ các
việc đã giao. Chúng tôi đề xuất \diva{}, một framework \emph{kiểm chứng chọn lọc} điều phối bằng
LangGraph, kết hợp ba lớp mô hình: (i) bộ cắt chủ đề DialSTART huấn luyện lại trên BERT tiếng Việt, không
cần biết trước số chủ đề; (ii) một \emph{mô hình quyết định} (decision model), loại mô hình mới nhận một
câu hỏi khai báo bằng ngôn ngữ tự nhiên cùng tập lựa chọn và trả về phân bố xác suất đã hiệu chỉnh thay vì
sinh văn bản, dùng làm lớp định tuyến; và (iii) các tác tử LLM sinh, trong đó tác tử kiểm chứng ReAct tra
cứu toàn bộ cuộc họp rồi đi tới đồng thuận với tác tử trích xuất. Thực nghiệm trên họp cơ quan nhà nước, họp doanh nghiệp/dự án và phiên họp Quốc hội cho thấy: bộ
chọn ranh giới nâng điểm cắt chủ đề trên dữ liệu thật từ 0,39 lên 0,64; mô hình quyết định nhận ra lời giao
việc tốt hơn hẳn luật từ khoá (F1 0,89 so với 0,68), với ECE 0,04, nên chỉ khoảng 20\% số lượt nói cần
đến bước kiểm chứng đắt tiền; và trên một cuộc họp có danh sách việc, hệ thống xác định đúng trạng thái
cả 4 việc.
\end{abstract}

%======================================================================
\section{Giới thiệu}
%======================================================================

Phần lớn quyết định trong một tổ chức được đưa ra trong cuộc họp: giao ban tuần của một UBND, họp chuyên
đề về một dự án đầu tư, họp giao ban của một công ty cấp nước, họp sprint của một nhóm phần mềm, hay họp
của ban quản lý dự án với nhà thầu. Bất kể bối cảnh, sau cuộc họp tổ chức cần một bản ghi lại \emph{ai
được giao làm gì, hạn khi nào, đã thống nhất điều gì}. Ở cơ quan nhà nước, đó là biên bản và thông báo
kết luận, một văn bản hành chính chính thức. Ở doanh nghiệp và nhóm dự án, đó là biên bản họp và danh
sách việc cần làm được đưa vào công cụ quản lý công việc. Ở cả hai nơi, danh sách việc này trở thành
chương trình kiểm điểm của cuộc họp tiếp theo.

Soạn biên bản bằng tay tốn công và dễ sót, nhất là với cuộc họp dài, nhiều nội dung. Nhận dạng giọng nói
tiếng Việt đã đủ tốt để tạo bản ghi, và mô hình ngôn ngữ lớn (LLM) viết tóm tắt rất trôi chảy; nhiều
công cụ họp trực tuyến cũng đã tự sinh tóm tắt và danh sách việc. Tuy nhiên, giữa \emph{một bản tóm tắt}
và \emph{một biên bản dùng được} có khoảng cách lớn. Sai sót trong biên bản dẫn tới việc bị bỏ sót, làm
trùng, hoặc bị quy sai trách nhiệm:
\begin{itemize}
  \item ghi một \emph{đề xuất} của một thành viên thành \emph{quyết định} của người có thẩm quyền;
  \item giao việc cho nhầm người (``anh Sơn'' khi cuộc họp có hai người tên Sơn), hoặc nhầm bên chủ trì
        với bên phối hợp;
  \item bỏ sót lời giao việc khẩu ngữ (``anh Phong gửi giúp em trong tuần này được không ạ?'') hay lời
        tự nhận việc (``bên em sẽ gửi trước thứ Sáu'');
  \item báo một việc ``đã xong'' trong khi người điều hành vừa yêu cầu làm lại.
\end{itemize}
Vì vậy người dùng cần kiểm tra được từng mục: mục này dựa vào câu nói nào, của ai.

\paragraph{Vì sao bài toán khó.} Thứ nhất, một cuộc họp có thể dài hàng trăm nghìn ký tự với nhiều nội
dung đan xen. Thứ hai, hành vi lời nói phụ thuộc vào vai người nói và văn phong: cùng câu ``đề nghị Sở
Xây dựng rà soát lại'' là chỉ đạo khi người chủ trì nói nhưng chỉ là kiến nghị khi đại biểu nói; trong họp
dự án, việc được giao hay tự nhận bằng khẩu ngữ, ít khi có từ ``giao''. Thứ ba, người nhận việc có thể là
cá nhân, phòng ban, đơn vị hay đối tác, nhiều bên với vai trò khác nhau, và cách gọi rất lỏng (``bên em'',
``team backend'', ``UBND TP''). Thứ tư, gần như chưa có dữ liệu tiếng Việt có nhãn cho bài toán này, dữ
liệu tổng hợp khác xa cuộc họp thật, và bộ tách người nói có thể gộp nhiều người vào một nhãn.

\paragraph{Hạn chế của các cách làm hiện có.} Gọi một LLM trên từng đoạn rồi tin kết quả thì không kiểm
soát được lỗi bịa và lỗi gán trách nhiệm. Cho nhiều tác tử tranh luận về mọi mục thì chi phí tăng theo
số mục, và mỗi cuộc tranh luận chỉ thấy đoạn hiện tại, trong khi lời chốt có thể nằm ở phần khác của
cuộc họp.

\paragraph{Cách tiếp cận.} \diva{} coi việc soạn biên bản là \emph{trích xuất có căn cứ} và tổ chức kiểm
chứng theo \emph{mức độ nghi ngờ}. Cuộc họp được cắt thành chủ đề và xử lý song song. Mỗi mục trích ra
phải chỉ được câu nói làm căn cứ. Một mô hình quyết định rẻ, đã hiệu chỉnh, phân loại hành vi lời nói và quyết định mục nào
chắc chắn và mục nào mơ hồ. Chỉ mục mơ hồ mới được chuyển cho một tác tử kiểm chứng, tác tử này tra cứu
toàn bộ cuộc họp và tranh luận có giới hạn với tác tử đã trích ra mục đó. Khi không thống nhất được hoặc
hệ thống gặp lỗi, mục được giữ lại kèm dấu hiệu cho người duyệt thay vì bị xoá. Kiến trúc không gắn với
một loại tổ chức: thông tin riêng của từng cuộc họp (người điều hành, danh sách thành viên và đơn vị,
danh sách việc cũ) là đầu vào tuỳ chọn.

\paragraph{Mô hình quyết định làm lớp định tuyến.} Điểm mới về công nghệ của \diva{} là dùng một
\emph{mô hình quyết định} (decision model) thay cho cả luật từ khoá lẫn việc nhờ LLM sinh ``trả lời một
nhãn''. Với mô hình quyết định, người thiết kế chỉ \emph{khai báo} câu hỏi: hướng dẫn và các lựa chọn, mỗi
lựa chọn kèm mô tả bằng ngôn ngữ tự nhiên (ví dụ ``giao việc: giao việc cho người/đơn vị cụ thể, hoặc một
người nhận làm việc''). Mô hình trả về phân bố xác suất trên các lựa chọn, không sinh văn bản tự do. Ba
tính chất của loại mô hình này quyết định thiết kế của \diva{}: xác suất được hiệu chỉnh nên ngưỡng có
nghĩa thực tế; cùng đầu vào cho cùng kết quả nên có thể lưu cache; và đủ nhanh (khoảng 0,3 giây mỗi lời
gọi) để chạy trên từng lượt nói. Định nghĩa nhãn mới chỉ cần viết thêm một câu hỏi, không cần huấn luyện.

\paragraph{Đóng góp.}
\begin{enumerate}
  \item Phát biểu bài toán sinh biên bản họp có căn cứ cho các cuộc họp ra quyết định, gồm bốn đầu ra và
        các yêu cầu chung: trích dẫn được, định danh người/đơn vị và vai trò, chuẩn hoá hạn chót, phân
        biệt quyết định với đề xuất (\S\ref{sec:problem}).
  \item \diva{}, một framework kiểm chứng chọn lọc trên LangGraph, trong đó một \emph{mô hình quyết
        định} đóng vai trò lớp định tuyến ở ba vị trí: cổng chọn tác tử, phân loại sơ bộ candidate, và
        kiểm trạng thái tiến độ; chỉ phần mơ hồ mới được chuyển sang cơ chế đồng thuận giữa tác tử trích
        xuất và tác tử kiểm chứng có công cụ tra cứu (\S\ref{sec:method}, \S\ref{sec:stack}).
  \item Một phương pháp cắt chủ đề cho cuộc họp tiếng Việt, không cần biết trước số chủ đề, giữ mỗi đoạn
        trong giới hạn xử lý của LLM (\S\ref{sec:seg}).
  \item Phân tích thực nghiệm trên ba bối cảnh họp (cơ quan nhà nước, doanh nghiệp/dự án, Quốc hội) cho
        thấy khoảng cách giữa dữ liệu tổng hợp và dữ liệu thật, giá trị của hiệu chỉnh xác suất khi định
        tuyến, và các điểm còn sai trong theo dõi tiến độ (\S\ref{sec:results}).
\end{enumerate}

%======================================================================
\section{Bài toán}
\label{sec:problem}
%======================================================================

\subsection{Bối cảnh họp}

Chúng tôi xét các cuộc họp có mục đích \emph{ra quyết định và phân công}, phân biệt với hội thảo hay buổi
trình bày một chiều. Ba bối cảnh tiêu biểu có đặc điểm khác nhau (Bảng~\ref{tab:settings}), nhưng cùng
cần một biên bản có căn cứ.

\begin{table}[t]
\centering\small
\begin{tabular}{@{}L{0.22\linewidth}L{0.7\linewidth}@{}}
\toprule
\textbf{Bối cảnh} & \textbf{Đặc điểm} \\
\midrule
Cơ quan nhà nước & Có người chủ trì rõ; chỉ đạo theo thứ bậc (``đề nghị Sở\ldots''); việc thường giao
cho đơn vị, có bên chủ trì và bên phối hợp; có thông báo kết luận. \\
Doanh nghiệp, dự án & Văn phong khẩu ngữ; nhiều việc tự nhận (``để em làm''); giao cho cá nhân, nhóm
hoặc đối tác; hạn tương đối (``thứ Sáu này'', ``sau sprint''). \\
Nghị trường, hội đồng & Phát biểu dài, một lượt chứa nhiều vấn đề; ít giao việc, nhiều kiến nghị và
biểu quyết. \\
\bottomrule
\end{tabular}
\caption{Ba bối cảnh họp và đặc điểm ảnh hưởng tới việc soạn biên bản.}
\label{tab:settings}
\end{table}

\subsection{Biên bản cần gì}

Bảng~\ref{tab:requirements} tóm tắt các yêu cầu chung, rút ra từ quy trình soạn biên bản và các biên bản
thật nhóm thu thập được. Cả bốn phần có chung một yêu cầu: \textbf{truy vết được}, tức mỗi mục chỉ ra
được câu nói làm căn cứ để người duyệt kiểm tra trong vài giây thay vì nghe lại cả cuộc họp.

\begin{table}[t]
\centering\small
\begin{tabular}{@{}L{0.2\linewidth}L{0.72\linewidth}@{}}
\toprule
\textbf{Phần} & \textbf{Yêu cầu} \\
\midrule
Diễn biến & Ý kiến chính của từng người theo nội dung; đúng người phát biểu. \\
Giao việc & Chỉ ghi khi có người có thẩm quyền giao hoặc có người tự nhận; đúng người/đơn vị và vai
trò chủ trì/phối hợp; hạn quy được về ngày. \\
Kết luận & Phân biệt nội dung đã chốt với đề xuất chưa được chấp nhận. \\
Tiến độ & Không bỏ việc cũ nào; biết việc nào chưa được báo cáo, quá hạn, bị yêu cầu làm lại. \\
\bottomrule
\end{tabular}
\caption{Bốn phần của biên bản và yêu cầu tương ứng.}
\label{tab:requirements}
\end{table}

\subsection{Phát biểu hình thức}

Đầu vào là bản ghi $\mathcal{T}=(u_1,\dots,u_n)$, mỗi lượt nói có người nói và nội dung. Kèm theo có thể
có ngày họp, người điều hành, danh sách người/đơn vị tham dự, và danh sách việc cần báo cáo $\mathcal{G}$
(mỗi việc có bên phụ trách, hạn và trạng thái trước cuộc họp). Đầu ra gồm bốn tập:
\begin{itemize}
  \item \textbf{Diễn biến} $\mathcal{D}$: luận điểm (người nói, nội dung, $E$);
  \item \textbf{Giao việc} $\mathcal{A}$: (nội dung, các bên nhận việc kèm vai trò, hạn, $E$, $t^{*}$);
  \item \textbf{Kết luận} $\mathcal{C}$: (nội dung, $E$, $t^{*}$);
  \item \textbf{Tiến độ} $\mathcal{P}$: mỗi $g\in\mathcal{G}$ có trạng thái mới trong \{xong, đang làm,
        vướng mắc, phải làm lại, hoãn, chỉ được nhắc, không được báo cáo\} kèm bằng chứng.
\end{itemize}
$E$ là tập lượt nói làm bằng chứng; $t^{*}$ là \emph{lượt quyết định} (lượt giao/nhận việc hoặc lượt
chốt). Ràng buộc căn cứ: mọi bằng chứng là lượt nói thật; mọi việc giao và kết luận sau kiểm chứng phải
có lượt quyết định, có thể nằm ở phần khác của cuộc họp; và không việc nào trong $\mathcal{G}$ bị bỏ
khỏi đầu ra.

\subsection{Thách thức}
\begin{description}
  \item[T1 Ngữ cảnh dài.] Phải cắt cuộc họp thành chủ đề, nhưng lời giao việc ở chủ đề này có thể được
        chốt ở chủ đề khác.
  \item[T2 Hành vi lời nói.] Giao việc, tự nhận việc, kết luận, đề xuất và báo cáo phân biệt theo
        \emph{nghĩa} và \emph{vai người nói}, không theo từ khoá; văn phong khác nhau giữa các bối cảnh.
  \item[T3 Trách nhiệm.] Người nhận việc là cá nhân, nhóm hay đơn vị, nhiều bên, cách gọi không chuẩn.
  \item[T4 Chi phí và độ tin cậy.] Kiểm chứng mọi mục bằng LLM quá đắt; không kiểm chứng thì lỗi bịa lọt ra.
  \item[T5 Dữ liệu ít và lệch.] Ít dữ liệu có nhãn, dữ liệu tổng hợp khác dữ liệu thật, nhãn người nói
        không tin cậy.
\end{description}

%======================================================================
\section{Công trình liên quan}
%======================================================================

\paragraph{Tóm tắt và trích xuất từ cuộc họp.} Các nghiên cứu trên AMI và ICSI
\citep{carletta2005ami,janin2003icsi} gồm phát hiện việc cần làm \citep{purver2007actionitems} và tóm tắt
theo truy vấn \citep{zhong2021qmsum}. Chúng chủ yếu cho tiếng Anh, ít đòi hỏi trích dẫn tới từng câu,
không định danh đơn vị và không theo dõi tiến độ việc qua các cuộc họp. Các công cụ họp thương mại sinh
tóm tắt và danh sách việc nhưng thường không cho biết mục nào dựa vào câu nói nào, và chưa hỗ trợ tốt
tiếng Việt.

\paragraph{Cắt chủ đề hội thoại.} Các phương pháp đi từ đo tương đồng giữa hai cửa sổ
\citep{hearst1997texttiling} và phân cụm trên embedding \citep{gklezakos2024treeseg} tới chấm điểm mạch
lạc giữa các lượt nói \citep{xing2021coherence,gao2023dialstart}. Hầu hết được đánh giá khi biết trước số
chủ đề; chọn độ mịn khi chạy thật là điểm chúng tôi tập trung xử lý.

\paragraph{Tác tử LLM và tự kiểm chứng.} ReAct \citep{yao2023react} xen kẽ suy luận và tra cứu. Tranh luận
đa tác tử \citep{du2023debate,liang2023divergent} và tự sửa lỗi \citep{madaan2023selfrefine,gou2024critic}
dùng phản hồi để cải thiện đầu ra. Định tuyến theo độ tin cậy
\citep{geifman2017selective,chen2023frugalgpt} chỉ chuyển sang mô hình đắt khi cần. \diva{} kết hợp hai
hướng: định tuyến bằng bộ phân loại đã hiệu chỉnh, và tranh luận \emph{có neo bằng chứng, có ngân sách,
có quy tắc an toàn}, chỉ dành cho phần mơ hồ.

\paragraph{Đầu ra có cấu trúc và mô hình quyết định.} Cách phổ biến để dùng LLM làm bộ phân loại là ép
đầu ra theo một JSON schema rồi đọc nhãn mô hình sinh ra. Cách này cho \emph{một} nhãn chứ không cho độ tin
cậy đáng tin; xác suất của mạng nơ-ron hiện đại vốn thường bị hiệu chỉnh kém \citep{guo2017calibration}.
Mô hình quyết định đi theo hướng khác: đầu ra là phân bố trên một tập lựa chọn được khai báo. \diva{} khai
thác đúng tính chất này để biến ngưỡng xác suất thành chính sách định tuyến giữa các tác tử.

%======================================================================
\section{Phương pháp \diva{}}
\label{sec:method}
%======================================================================

\begin{figure*}[t]
\centering
\begin{tikzpicture}[
  font=\small, x=0.96cm,
  box/.style={draw, rounded corners=2pt, align=center, minimum height=2.6em, inner sep=3pt, fill=white},
  dec/.style={draw, diamond, aspect=2, align=center, inner sep=0pt, fill=white},
  arr/.style={-{Latex[length=2mm]}}]
  \node[box] (asr) at (-0.4,0) {Bản ghi\\giọng nói};
  \node[box] (seg) at (1.6,0) {Cắt chủ đề\\(\S\ref{sec:seg})\\{\scriptsize\itshape\color{blue!50!black} PyTorch, ViBERT}};
  \node[box] (ext) at (5.6,0) {Trích xuất theo chủ đề\\diễn biến $\cdot$ giao việc\\kết luận $\cdot$ tiến độ\\{\scriptsize\itshape\color{blue!50!black} LLM sinh (gemma)}};
  \node[dec] (tri) at (8.9,0) {Mô hình\\quyết định\\{\scriptsize\itshape\color{blue!50!black} Decisions API}};
  \node[box] (ok)  at (11.6,1.15) {Chấp nhận};
  \node[box] (ver) at (11.6,-1.15) {Kiểm chứng tra cứu\\+ đồng thuận\\{\scriptsize\itshape\color{blue!50!black} ReAct, GPT-4o-mini}};
  \node[box] (out) at (14.2,0) {Biên bản\\+ vết\\kiểm chứng\\{\scriptsize\itshape\color{blue!50!black} FastAPI, Redis}};
  \node[box, fill=gray!10, minimum height=1.8em] (reg) at (5.6,2.75)
       {Danh bạ toàn cuộc họp: người nói, người điều hành, đơn vị, việc cũ};
  \draw[arr] (asr) -- (seg);
  \draw[arr] (seg) -- node[below, font=\scriptsize]{$K$ chủ đề} (ext);
  \draw[arr] (ext) -- (tri);
  \draw[arr] (tri.north) |- node[pos=0.25, left, font=\scriptsize]{chắc chắn} (ok.west);
  \draw[arr] (tri.south) |- node[pos=0.25, left, font=\scriptsize]{mơ hồ} (ver.west);
  \draw[arr] (ok.east) -| (out.north);
  \draw[arr] (ver.east) -| (out.south);
  \draw[arr, dashed] (reg) -- (ext);
  \node[above=1pt of reg, font=\scriptsize\itshape, text=blue!50!black] {điều phối map--reduce bằng LangGraph};
  \begin{scope}[on background layer]
    \node[draw, dashed, rounded corners, fill=blue!4, inner xsep=5pt, inner ysep=4pt,
          fit=(ext) (tri) (ok) (ver),
          label={[font=\scriptsize]below:xử lý song song cho từng chủ đề}] {};
  \end{scope}
\end{tikzpicture}
\caption{Tổng quan \diva{}. Kiểm tra rẻ chạy cho mọi mục; chỉ các mục mơ hồ được chuyển sang tác tử
kiểm chứng có công cụ tra cứu toàn cuộc họp và cơ chế đồng thuận.}
\label{fig:overview}
\end{figure*}

Hình~\ref{fig:overview} tóm tắt \diva{}. Thiết kế dựa trên ba nguyên tắc:
\begin{description}
  \item[P1 Căn cứ trước, nội dung sau.] Mọi tác tử phải chỉ ra lượt nói làm căn cứ trước khi viết nội
        dung hay đưa ra kết luận.
  \item[P2 Chi phí tỉ lệ với mức nghi ngờ.] Kiểm tra rẻ chạy cho mọi mục; kiểm chứng đắt chỉ chạy cho
        mục mơ hồ.
  \item[P3 An toàn khi không chắc.] Lỗi hệ thống hay bất đồng không giải quyết được thì \emph{giữ} mục
        kèm dấu hiệu cho người duyệt, không tự xoá.
\end{description}

\subsection{Cắt chủ đề dựa trên mạch lạc}
\label{sec:seg}

Chúng tôi huấn luyện lại DialSTART \citep{gao2023dialstart} cho tiếng Việt. Với mỗi khe $g$ giữa hai lượt
nói liên tiếp, mô hình cho điểm liền mạch kết hợp hai tín hiệu: lượt sau có \emph{nối tiếp} lượt trước
không, và hai cửa sổ hai bên có \emph{cùng chủ đề} không:
\begin{equation}
\mathrm{score}(g)=\sigma\big(\mathrm{coh}(u_{\le g},u_{g+1})+\cos(\bar{\mathbf{e}}_{\text{trái}},\bar{\mathbf{e}}_{\text{phải}})\big).
\end{equation}
Khe có điểm ``lõm'' sâu so với hai đỉnh lân cận là ứng viên ranh giới.

Mô hình gốc được đánh giá với \emph{số chủ đề đã biết}, điều không có khi chạy thật. Chúng tôi đề xuất bộ
chọn ranh giới bốn bước: (1) làm mượt điểm để bỏ các chỗ lõm giả do câu xen ngang (``Vâng'', ``Mời
anh\ldots''); (2) giữ khe sâu hơn một ngưỡng thích nghi theo phân bố điểm của chính cuộc họp đó; (3) bỏ
đoạn quá ngắn, với độ dài tối thiểu tỉ lệ theo độ dài cuộc họp; (4) tách đoạn vượt giới hạn xử lý của LLM.
Nhờ vậy cùng một cấu hình chạy được cho cuộc họp từ 5 nghìn tới 260 nghìn ký tự.

\subsection{Xử lý song song với ngữ cảnh toàn cuộc họp}

Xử lý chủ đề tuần tự để mang thông tin người nói sang chủ đề sau làm thời gian tăng tuyến tính theo số
chủ đề. \diva{} dựng \emph{một lần} một \textbf{danh bạ} cho cả cuộc họp (người nói, người điều hành,
người/đơn vị tham dự, danh sách việc), rồi xử lý mọi chủ đề song song và độc lập. Mỗi chủ đề có các tác
tử trích xuất riêng cho diễn biến, giao việc, kết luận và tiến độ. Prompt mã hoá các quy tắc chung của
cuộc họp ra quyết định, ví dụ: \emph{lời ``đề nghị/yêu cầu + người/đơn vị + việc'' của người có thẩm quyền
là giao việc, kể cả khi không có chữ ``giao''; cùng câu đó của một thành viên khác chỉ là đề xuất}; và
\emph{lời một người nói mình sẽ làm việc gì mà không ai phản đối là tự nhận việc}. Kết quả được ghép theo
thứ tự chủ đề và gộp mục trùng giữa các chủ đề.

\subsection{Mô hình quyết định}
\label{sec:decision}

Một mô hình quyết định $\mathcal{M}$ nhận đầu vào văn bản $x$ và một \emph{câu hỏi khai báo}
$q=(\text{hướng dẫn},\{(c_j,\text{mô tả}_j)\}_{j=1}^{m})$, rồi trả về phân bố
$\mathcal{M}(x,q)\in\Delta^{m-1}$ trên các lựa chọn $c_1,\dots,c_m$. Khác với LLM sinh, mô hình không
sinh văn bản, nên không cần phân tích đầu ra và không có lỗi định dạng. \diva{} dùng bốn câu hỏi
(Bảng~\ref{tab:questions}), đều viết bằng tiếng Việt, với đầu vào là lượt cần xét kèm 2 lượt ngữ cảnh
đứng trước. Lượt dài hơn 2.500 ký tự được cắt theo câu thành nhiều cửa sổ, cửa sổ trước làm ngữ cảnh cho
cửa sổ sau. Kết quả được nhớ theo nội dung đầu vào, nên cổng chọn tác tử và bước phân loại sơ bộ dùng
chung lời gọi khi xét cùng một lượt.

\begin{table}[t]
\centering\small
\begin{tabular}{@{}L{0.24\linewidth}L{0.68\linewidth}@{}}
\toprule
\textbf{Câu hỏi} & \textbf{Lựa chọn và vai trò} \\
\midrule
Hành vi lời nói & giao việc, kết luận, đề xuất, bác bỏ, thảo luận, khác. Cổng chọn tác tử và phân loại
sơ bộ (\S\ref{sec:routing}). \\
Kiểu việc & báo cáo tại họp / làm ra sản phẩm. Chọn câu hỏi kiểm tiến độ (\S\ref{sec:progress}). \\
Trạng thái việc & báo cáo xong, đang làm, vướng mắc, yêu cầu làm lại, hoãn, chỉ nhắc, không liên quan. \\
Trạng thái báo cáo & đã báo cáo, chưa có báo cáo, xin hoãn, yêu cầu báo cáo lại, chỉ nhắc, không liên quan. \\
\bottomrule
\end{tabular}
\caption{Bốn câu hỏi khai báo cho mô hình quyết định.}
\label{tab:questions}
\end{table}

\subsection{Định tuyến theo hành vi lời nói}
\label{sec:routing}

Gọi $h(\cdot\mid u)=\mathcal{M}(u,q_{\text{hành vi}})$ là phân bố trên sáu hành vi của lượt $u$: \emph{giao
việc} (gồm cả tự nhận việc), \emph{kết luận}, \emph{đề xuất}, \emph{bác bỏ}, \emph{thảo luận}, \emph{khác}.
Hướng dẫn của câu hỏi yêu cầu xét theo nghĩa và theo vai người nói. Xác suất ``chốt'' là
\begin{equation}
\pchot(u)=h(\text{giao việc}\mid u)+h(\text{kết luận}\mid u).
\end{equation}
Vì xác suất này được hiệu chỉnh tốt (\S\ref{sec:rq2}), hai ngưỡng $\tau_{\text{thấp}}=0{,}2$ và
$\tau_{\text{cao}}=0{,}8$ chia lượt nói thành ba vùng có ý nghĩa thực tế: \emph{rõ là không chốt},
\emph{mơ hồ}, \emph{rõ là chốt}. Câu hỏi này được dùng ở hai chỗ:
\begin{itemize}
  \item \textbf{Cổng trích xuất.} Tác tử giao việc và kết luận chỉ chạy trên chủ đề có ít nhất một lượt
        với $\pchot\ge\tau_{\text{thấp}}$. Ngưỡng thấp có chủ đích: chạy thừa chỉ tốn chi phí, bỏ nhầm
        thì mất nội dung.
  \item \textbf{Phân loại sơ bộ.} Lượt quyết định mà tác tử khai cho mỗi mục được đưa qua $h$. Vùng
        \emph{rõ là chốt} được chấp nhận ngay nếu các kiểm tra luật khác cũng đạt (người nhận việc có
        thật, không mơ hồ giữa nhiều người); các vùng còn lại chuyển sang kiểm chứng.
\end{itemize}

\subsection{Kiểm chứng có tra cứu và đồng thuận}
\label{sec:consensus}

\paragraph{Tác tử kiểm chứng} hoạt động theo kiểu ReAct \citep{yao2023react} với ba công cụ chỉ đọc trên
\emph{toàn bộ cuộc họp}: đọc một lượt nói, tìm lượt liên quan theo từ khoá, và tra ứng viên người/đơn vị
cho một cách gọi. Kết luận là \emph{giữ}, \emph{sửa}, \emph{bỏ} hoặc \emph{chưa kết luận được}. Tác tử
phải chỉ ra lượt quyết định \emph{trước} khi kết luận (P1); kết luận giữ/sửa mà không có lượt quyết định
có thật bị hạ thành bỏ. Dấu hiệu nghi ngờ của bước phân loại sơ bộ chỉ là gợi ý, không phải bằng chứng,
và số bước tra cứu bị giới hạn.

\paragraph{Đồng thuận.} Khi tác tử kiểm chứng không giữ nguyên một mục, phản hồi được gửi lại cho chính
tác tử đã trích ra mục đó. Tác tử này chọn \emph{chấp nhận}, \emph{sửa} (kèm lượt quyết định có thật) hoặc
\emph{bảo vệ} (chỉ ra lượt quyết định bị bỏ sót). Hai bên trao đổi tối đa $R$ vòng; quy tắc kết thúc ở
Bảng~\ref{tab:termination} cụ thể hoá P3. Mọi vòng được lưu thành \emph{vết kiểm chứng} để người duyệt
xem vì sao một mục được giữ, sửa hay bỏ.

\begin{table}[t]
\centering\small
\begin{tabular}{@{}L{0.56\linewidth}L{0.36\linewidth}@{}}
\toprule
\textbf{Kết thúc} & \textbf{Kết quả} \\
\midrule
Hai bên thống nhất & Theo kết quả thống nhất \\
Hết vòng, bên kiểm chứng đã kết luận & Theo bên kiểm chứng \\
Hết vòng, vẫn chưa kết luận được & Giữ mục \textbf{gốc}, đánh dấu \\
Lỗi hệ thống & Giữ mục \textbf{gốc}, đánh dấu \\
\bottomrule
\end{tabular}
\caption{Quy tắc kết thúc của cơ chế đồng thuận.}
\label{tab:termination}
\end{table}

\subsection{Định danh trách nhiệm}

Mỗi cách gọi người nhận việc (``anh Sơn'', ``team backend'', ``Sở Tài chính chủ trì, phối hợp Sở Xây
dựng'') được tách thành các bên kèm vai trò (chủ trì, phối hợp, cùng thực hiện). Mỗi bên được chấm điểm
ứng viên theo ba nguồn, lần lượt ưu tiên: \emph{tên} (khớp đầy đủ, một phần, hay chỉ tên gọi); \emph{ngữ
cảnh hội thoại} (người đáp lời ngay sau lời giao thường là người nhận việc; người giao việc không tự gọi
mình bằng tên); và \emph{chức năng} (chức danh hay chức năng của đơn vị khớp nội dung việc). Khi không ứng
viên nào đủ rõ, việc giao vẫn được giữ kèm danh sách ứng viên và điểm để người duyệt chọn; hệ thống không
tự đoán.

\subsection{Theo dõi tiến độ việc đã giao}
\label{sec:progress}

Khi có danh sách việc cũ, tác tử tiến độ trích trạng thái từng việc được nhắc tới, rồi trạng thái được
kiểm bằng mô hình quyết định. Quan sát then chốt là có \textbf{hai kiểu việc}: \emph{việc báo cáo} (``báo cáo
tình hình các dự án'', ``cập nhật số liệu doanh thu'') xong khi nội dung đã được trình bày tại cuộc họp,
còn \emph{việc làm ra sản phẩm} (tờ trình, hồ sơ, bản thiết kế, tính năng phần mềm) xong khi sản phẩm đã
có. Với việc báo cáo, tình trạng của \emph{đối tượng} được báo cáo (``dự án đang làm thủ tục'') là nội
dung báo cáo, không phải trạng thái của \emph{việc}. Vì vậy \diva{} phân loại kiểu việc trước rồi đặt câu
hỏi kiểm phù hợp.

Các cập nhật của cùng một việc được gộp theo hai quy tắc: \textbf{lời chỉ đạo của người có thẩm quyền
thắng lời báo cáo} (báo cáo ``xong'' rồi người điều hành yêu cầu làm lại thì kết quả là \emph{phải làm
lại}, kèm cờ mâu thuẫn), và cùng loại thì lời muộn hơn thắng. Việc không ai nhắc tới được đánh dấu
\emph{không được báo cáo}; việc quá hạn mà chưa xong được đánh dấu \emph{quá hạn}. Mỗi việc giao mới cũng
được xếp quan hệ với danh sách cũ: \emph{việc mới}, \emph{tiếp nối}, \emph{giao lại cho bên khác} hoặc
\emph{đổi hạn}, giúp cập nhật sổ theo dõi hay công cụ quản lý công việc.

%======================================================================
\section{Hiện thực hoá và tech stack}
\label{sec:stack}
%======================================================================

\diva{} được hiện thực thành một dịch vụ hoàn chỉnh. Bảng~\ref{tab:stack} liệt kê công nghệ ở từng lớp.

\begin{table}[t]
\centering\small
\begin{tabular}{@{}L{0.26\linewidth}L{0.66\linewidth}@{}}
\toprule
\textbf{Lớp} & \textbf{Công nghệ} \\
\midrule
Cắt chủ đề & PyTorch, HuggingFace Transformers; DialSTART trên \code{vibert4news-base-cased} \\
Điều phối tác tử & LangGraph: \code{StateGraph}, subgraph theo chủ đề, \code{Send} cho map--reduce, reducer cộng dồn \\
LLM sinh & gemma tự triển khai (endpoint tương thích OpenAI) cho gán nhãn chủ đề và trích xuất; GPT-4o-mini với JSON schema strict cho tác tử kiểm chứng \\
Mô hình quyết định & OpenAI Decisions API (\code{gpt-6-luna}) \\
Công cụ tra cứu & BM25 trên âm tiết và cặp âm tiết tiếng Việt, không cần CSDL vector \\
Dịch vụ & FastAPI; Redis cho cache kết quả và khoá chạy-một-lần; Docker Compose \\
Giao diện & React 19, TypeScript, Vite, Tailwind CSS \\
\bottomrule
\end{tabular}
\caption{Tech stack của \diva{}.}
\label{tab:stack}
\end{table}

\paragraph{Điều phối bằng LangGraph.} Pipeline là một đồ thị trạng thái hai tầng. Đồ thị cha dựng danh
bạ cuộc họp rồi phát mỗi chủ đề vào một subgraph bằng \code{Send}; các subgraph chạy song song và trả
kết quả qua reducer cộng dồn; nút tổng hợp sắp lại theo thứ tự chủ đề. Trong subgraph, các tác tử trích
xuất là các nhánh song song, và mỗi candidate mơ hồ lại được \code{Send} riêng sang tác tử kiểm chứng.
Một semaphore dùng chung bọc mọi lời gọi LLM để giới hạn tải bất kể số chủ đề.

\paragraph{ReAct không cần tool calling.} Tác tử kiểm chứng được dựng trên giao diện sinh JSON một lượt:
mỗi bước mô hình trả \{suy nghĩ, hành động, tham số, lượt quyết định, kết luận\}, mã nguồn chạy công cụ và
đưa observation vào bước sau. Nhờ vậy tác tử chạy được trên mọi nhà cung cấp tương thích OpenAI, kể cả
mô hình mở không hỗ trợ tool calling. Thứ tự trường trong schema được dùng như một ràng buộc suy luận:
lượt quyết định đứng trước kết luận.

\paragraph{Phân vai mô hình.} Mỗi loại mô hình làm đúng việc nó giỏi: mô hình mã hoá chuyên biệt cho cắt
chủ đề (nhanh, chạy cục bộ), mô hình quyết định cho các câu hỏi phân loại cần xác suất, LLM mở tự triển
khai cho phần trích xuất chiếm nhiều token, và một LLM khác họ cho bước kiểm chứng để giảm thiên lệch tự
đồng ý.

\paragraph{Vận hành.} Kết quả được cache theo hash của nội dung yêu cầu cùng dấu vân tay của mã nguồn và
cấu hình, nên đổi prompt là tự vô hiệu cache cũ; từng phần biên bản có thể chạy lại riêng mà không cắt lại
chủ đề. Giao diện cho người duyệt bấm vào từng mục để xem câu nói làm căn cứ và vết kiểm chứng.

%======================================================================
\section{Thiết lập thực nghiệm}
%======================================================================

\paragraph{Câu hỏi nghiên cứu.}
\textbf{RQ1}: Có cắt được chủ đề cuộc họp thật mà không cần biết trước số chủ đề?
\textbf{RQ2}: Mô hình quyết định có định tuyến kiểm chứng tốt hơn luật từ khoá, trên các bối cảnh
họp khác nhau?
\textbf{RQ3}: Có thể bỏ qua trích xuất ở chủ đề không có giao việc/kết luận mà không mất nội dung?
\textbf{RQ4}: Theo dõi tiến độ có đúng không, và yếu tố nào quyết định độ chính xác?
\textbf{RQ5}: Thời gian chạy trên cuộc họp dài có chấp nhận được?

\paragraph{Dữ liệu.} Bảng~\ref{tab:data} liệt kê các tập dữ liệu. Tập tổng hợp được sinh có chủ đích để
phủ cả ba bối cảnh: khoảng 12 cuộc họp cơ quan nhà nước (UBND, sở ngành, công an, phòng chống thiên tai),
14 cuộc họp doanh nghiệp và dự án (phần mềm, xây dựng, cấp nước, ngân hàng, bán lẻ, bệnh viện, trường
học), và 4 cuộc họp chỉ có nhãn người nói ẩn danh như đầu ra của bộ tách người nói.

\begin{table}[t]
\centering\small
\begin{tabular}{@{}L{0.27\linewidth}L{0.65\linewidth}@{}}
\toprule
\textbf{Tập} & \textbf{Mô tả} \\
\midrule
Quốc hội & 47 phiên Quốc hội khoá XIII có chủ đề gán tay; 12 phiên dùng đánh giá cắt chủ đề \\
Họp thật (gold) & 3 cuộc họp cơ quan nhà nước, ranh giới chủ đề đã duyệt \\
Tổng hợp & 30 cuộc họp thuộc ba bối cảnh; 869 lượt có nhãn hành vi lời nói; 166 kết luận gold \\
Có danh sách việc & 1 cuộc họp thật (32 lượt, 3 chủ đề), 4 việc có trạng thái gold \\
Họp dài & 1 cuộc họp thật, 1.259 lượt nói, 14 chủ đề \\
\bottomrule
\end{tabular}
\caption{Dữ liệu thực nghiệm.}
\label{tab:data}
\end{table}

\paragraph{Độ đo.} Cắt chủ đề: $P_k$ \citep{beeferman1999pk}, F1 ranh giới theo dung sai 2k/5k ký tự, và
điểm tổng hợp $S=\tfrac14(\mathrm{F1}_{2k}+\mathrm{F1}_{5k}+2-2P_k)$. Hành vi lời nói: P, R, F1 và sai số
hiệu chỉnh ECE \citep{guo2017calibration}. Tiến độ: số việc có trạng thái khớp gold. Thời gian chạy.

\paragraph{Mô hình.} Theo Bảng~\ref{tab:stack}. Các so sánh ở RQ2--RQ4 dùng cùng mô hình quyết định với
các câu hỏi ở Bảng~\ref{tab:questions}; đường cơ sở là bộ luật từ khoá của hệ thống trước. $R=3$ vòng
đồng thuận, tối đa 3 lần tra cứu mỗi vòng.

%======================================================================
\section{Kết quả và phân tích}
\label{sec:results}
%======================================================================

Các kết quả dưới đây đánh giá từng thành phần trên các tập khác nhau; chất lượng end-to-end của phần
Giao việc và Kết luận chưa được đo trên dữ liệu thật có nhãn (xem phần Hạn chế).

\subsection{RQ1: Cắt chủ đề}

\begin{table}[t]
\centering\small
\begin{tabular}{@{}lll@{}}
\toprule
\textbf{Cấu hình} & \textbf{Dữ liệu} & \textbf{Kết quả} \\
\midrule
Biết trước số chủ đề & Tổng hợp & $P_k=0{,}011$ \\
Chưa có bộ chọn & Thật & $P_k=0{,}485$ \\
+ bộ chọn (ban đầu) & Thật / TH & $S=0{,}392$ / $0{,}336$ \\
\textbf{+ bộ chọn (đề xuất)} & Thật / TH & $S=\mathbf{0{,}637}$ / $\mathbf{0{,}686}$ \\
\bottomrule
\end{tabular}
\caption{Khoảng cách giữa dữ liệu tổng hợp (TH) và dữ liệu thật trong cắt chủ đề.}
\label{tab:seg}
\end{table}

Bảng~\ref{tab:seg} cho thấy mô hình gần như hoàn hảo trên dữ liệu tổng hợp nhưng hỏng trên cuộc họp thật
(cắt quá vụn). Bộ chọn ranh giới thu hẹp phần lớn khoảng cách. Tương quan giữa điểm trên hai loại dữ liệu
chỉ ở mức vừa (Spearman 0,56), nên không nên tinh chỉnh chỉ dựa trên dữ liệu tổng hợp. Trên 12 phiên Quốc
hội, bỏ bước làm mượt làm F1 giảm từ 0,418 xuống 0,297, là bước có ảnh hưởng lớn nhất. Bối cảnh nghị
trường cũng cho thấy giới hạn của việc cắt theo lượt nói: ít nhất 49,5\% cặp (đại biểu, phiên họp) có một
lượt phát biểu chứa nhiều chủ đề, nên mọi phương pháp cắt theo lượt nói đều có một mức sai số không tránh
được ở bối cảnh này.

\subsection{RQ2: Hành vi lời nói và định tuyến kiểm chứng}
\label{sec:rq2}

\begin{table}[t]
\centering\small
\begin{tabular}{@{}lccc@{}}
\toprule
\textbf{Phương pháp} & \textbf{P} & \textbf{R} & \textbf{F1} \\
\midrule
Luật từ khoá & 0,774 & 0,610 & 0,682 \\
Mô hình quyết định, $p\ge0{,}5$ & \textbf{0,863} & \textbf{0,928} & \textbf{0,894} \\
Mô hình quyết định, $p\ge0{,}8$ & 0,923 & 0,816 & 0,866 \\
\bottomrule
\end{tabular}
\caption{Nhận diện lượt chốt (giao việc/kết luận) trên 869 lượt nói của tập tổng hợp.}
\label{tab:turnact}
\end{table}

\begin{table}[t]
\centering\small
\begin{tabular}{@{}lccl@{}}
\toprule
\textbf{Vùng} & \textbf{\% lượt} & \textbf{\% là chốt} & \textbf{Xử lý} \\
\midrule
$p<0{,}2$ & 31 & 4 & không chốt \\
$0{,}2\le p<0{,}8$ & 20 & 45 & kiểm chứng \\
$p\ge0{,}8$ & 49 & 92 & chấp nhận \\
\bottomrule
\end{tabular}
\caption{Ba vùng định tuyến theo $\pchot$.}
\label{tab:zones}
\end{table}

Mô hình quyết định vượt luật từ khoá rõ rệt (Bảng~\ref{tab:turnact}), với ECE 0,040: khi mô hình nói 90\%
thì thực tế đúng xấp xỉ 90\%. Tính chất này cho phép dùng ngưỡng như một chính sách vận hành.
Bảng~\ref{tab:zones} cho thấy chỉ khoảng 20\% lượt nói thật sự mơ hồ, với tỉ lệ đúng gần như tung đồng
xu; đây đúng là phần nên dành cho bước kiểm chứng đắt tiền.

Lợi ích lớn nhất nằm ở \textbf{lời giao việc khẩu ngữ}, vốn phổ biến trong họp doanh nghiệp và dự án:
tỉ lệ nhận ra tăng từ 51\% lên 91\%. Câu ``Anh Phong gửi giúp em trong tuần này được không ạ?'' bị luật
từ khoá bỏ qua vì là câu hỏi, còn mô hình quyết định cho $p=0{,}95$; lời đáp ``Được, thứ Sáu này anh gửi cả
hai'' cũng được nhận là tự nhận việc ($p=0{,}99$). Hai loại còn khó là \emph{việc chưa có người nhận} và
\emph{quyết định bị thay thế sau đó}; cả hai cần ngữ cảnh rộng hơn vài lượt, đó là lý do tác tử kiểm chứng
được tra cứu toàn bộ cuộc họp.

\subsection{RQ3: Bỏ qua trích xuất an toàn}

\begin{table}[t]
\centering\small
\begin{tabular}{@{}L{0.52\linewidth}L{0.4\linewidth}@{}}
\toprule
\textbf{Tiêu chí bỏ qua} & \textbf{Nội dung gold bị mất} \\
\midrule
Không có từ khoá giao/chốt & 1 việc giao + 2 kết luận (chỉ bỏ được 3/304 lần chạy) \\
Lọc riêng theo $p(\text{kết luận})$ & 5/166 kết luận \\
\textbf{Không lượt nào có $\pchot\ge0{,}2$} & \textbf{0} \\
\bottomrule
\end{tabular}
\caption{Mô phỏng các tiêu chí bỏ qua tác tử giao việc/kết luận trên 30 cuộc họp tổng hợp.}
\label{tab:gate}
\end{table}

Từ khoá vừa không tiết kiệm được gì vừa làm mất nội dung (Bảng~\ref{tab:gate}). Lọc riêng theo từng nhãn
cũng mất kết luận, vì một số lời chốt kết luận bị xếp vào nhãn giao việc. Gộp hai nhãn thành $\pchot$ ở
ngưỡng thấp thì không mất mục nào. Phần tiết kiệm bị chặn trên bởi tỉ lệ chủ đề không có giao việc lẫn
kết luận, khoảng 18\% trên dữ liệu tổng hợp.

\subsection{RQ4: Theo dõi tiến độ}

Trên cuộc họp có danh sách việc (3 việc thật giao từ trước và 1 việc giả lập không được nhắc tới),
\diva{} xác định đúng trạng thái \textbf{4/4 việc} ở cả 4 lần chạy. Việc giả lập được đánh dấu đúng là
\emph{không được báo cáo}; yêu cầu ``làm lại tờ trình'' được xếp là \emph{tiếp nối} việc cũ, còn các việc
phát sinh trong cuộc họp là \emph{việc mới}. Cả cuộc họp tốn 20 lời gọi mô hình quyết định, so với 13
khi không có danh sách việc. Ba yếu tố quyết định độ chính xác:

\paragraph{Chất lượng danh sách việc đầu vào.} Ở phiên bản danh sách đầu tiên, 3/6 ``việc cũ'' thực ra
phát sinh ngay trong cuộc họp. Hệ thống chỉ khớp 3/6, và mỗi lần chạy khớp một bộ 3 việc khác nhau. Định
nghĩa thế nào là ``việc cũ'' là quyết định của người dùng, không phải lỗi mô hình.

\paragraph{Câu hỏi kiểm phải hợp với kiểu việc.} Với việc ``báo cáo tình hình các dự án'', câu hỏi chung
xếp lượt trả lời là ``đang làm'' (0,92) vì các dự án đang làm thủ tục; câu hỏi dành cho việc báo cáo xếp
đúng là ``đã báo cáo'' (0,99). Câu hỏi kiểu việc tách được việc báo cáo (1,00) khỏi việc làm ra sản
phẩm (tờ trình 0,35; tiếp thu góp ý 0,03).

\paragraph{Phân công giữa LLM sinh và mô hình quyết định.} Câu ``đã giải trình, tiếp thu'' nằm trong một đoạn dài
nói về đề án ``tiếp tục hoàn thiện'' bị mô hình quyết định xếp ``đang làm'', và ba cách viết lại câu hỏi đều
không sửa được mà không làm hỏng trường hợp khác. Vì vậy tác tử, vốn đọc cả chủ đề, quyết ``xong hay đang
làm''; mô hình quyết định chỉ quyết lượt nói là báo cáo, chỉ đạo hay chỉ nhắc tới, và ngăn một kiến nghị của
thành viên biến việc đã xong thành ``phải làm lại''.

Ngoài ra, LLM mở vẫn có thể cho kết quả khác nhau giữa các lần chạy với cùng đầu vào, và cắt chủ đề khác
đi cũng làm prompt thay đổi; so sánh hai cấu hình chỉ có ý nghĩa khi cố định phân đoạn và chạy lặp.

\subsection{RQ5: Thời gian chạy}

Cuộc họp 1.259 lượt, 14 chủ đề được xử lý trong khoảng \textbf{78 giây}, trong đó cắt chủ đề mất 12 giây
(đo trước khi có vòng đồng thuận). Nhờ xử lý song song, thời gian không còn tăng tuyến tính theo số chủ
đề. Chi phí kiểm chứng mỗi mục mơ hồ bị chặn trên ở 15 lần gọi LLM; trong ví dụ điển hình, tác tử kiểm
chứng tìm thấy lời giao việc thật của người điều hành ở lượt khác, tác tử trích xuất sửa người nhận việc
theo đó, và hai bên thống nhất sau 2 vòng, tổng cộng 4 lần gọi.

%======================================================================
\section{Thảo luận}
%======================================================================

\paragraph{Ý nghĩa thực tiễn.} \diva{} không thay người soạn biên bản mà chuyển công việc từ \emph{nghe
lại và viết} sang \emph{duyệt có trích dẫn}. Mỗi mục kèm câu nói làm căn cứ; mục mơ hồ được đánh dấu kèm
lý do và ứng viên người nhận việc; vết kiểm chứng cho biết vì sao một mục được giữ hay bỏ. Thiết kế
``không chắc thì giữ lại và đánh dấu'' chấp nhận nhiều mục cần duyệt hơn để đổi lấy việc không âm thầm mất
nội dung; ở mọi bối cảnh, sót một việc được giao thường đắt hơn duyệt thừa một mục.

\paragraph{Tổng quát hoá giữa các bối cảnh.} Phần lớn khác biệt giữa các bối cảnh nằm ở \emph{đầu vào
tuỳ chọn} (người điều hành, danh sách thành viên và đơn vị, danh sách việc cũ) và ở văn phong, chứ không ở
kiến trúc. Câu hỏi hành vi lời nói của mô hình quyết định xét theo nghĩa và vai người nói nên xử lý được cả chỉ đạo theo thứ
bậc lẫn tự nhận việc khẩu ngữ. Những điểm còn phụ thuộc bối cảnh là quy tắc ``ai có thẩm quyền giao việc''
(rõ ở cơ quan nhà nước, mờ hơn ở nhóm dự án) và cách đặt tên đơn vị.

\paragraph{Giao thức đánh giá đề xuất.} Bước tiếp theo là xây gold end-to-end gồm biên bản thật của
\emph{cả ba} bối cảnh, có danh sách việc, hai người gán nhãn độc lập, và đo các ablation: tắt cổng trích
xuất; danh bạ toàn cuộc họp so với ngữ cảnh chuyển dồn; thay đổi hoặc bỏ vòng đồng thuận; kiểm chứng tra
cứu toàn cuộc họp so với tranh luận trong đoạn; các bộ cắt chủ đề khác nhau; câu hỏi tiến độ chung so với
tuỳ kiểu việc. Kết quả nên được báo cáo riêng theo từng bối cảnh.

%======================================================================
\section{Kết luận}
%======================================================================

Chúng tôi phát biểu bài toán sinh biên bản họp tiếng Việt có căn cứ cho các cuộc họp ra quyết định và
phân công, trong đó đúng \emph{trách nhiệm} quan trọng hơn đẹp câu chữ, và đề xuất \diva{}. \diva{} cắt
cuộc họp thành chủ đề, xử lý song song bằng LangGraph, và dùng một mô hình quyết định đã hiệu chỉnh để quyết
định nơi cần trích xuất và nơi cần kiểm chứng; việc kiểm chứng dành cho một tác tử tra cứu toàn cuộc họp,
làm việc theo cơ chế đồng thuận có quy tắc an toàn. Thực nghiệm thành phần trên ba bối cảnh họp cho thấy
hiệu chỉnh xác suất là điều kiện để định tuyến hiệu quả, từ khoá không đủ để nhận ra lời giao việc tiếng
Việt, và theo dõi tiến độ cần phân biệt việc báo cáo với việc làm ra sản phẩm. Đánh giá end-to-end trên dữ
liệu thật mở rộng là bước tiếp theo cần thiết.

%======================================================================
\section*{Hạn chế (Limitations)}
%======================================================================

\textbf{Chưa có đánh giá end-to-end} cho Giao việc và Kết luận trên dữ liệu thật có nhãn, và chưa kiểm
chứng được rằng vòng đồng thuận không khiến một bên nhượng bộ sai. \textbf{Dữ liệu thật lệch về cơ quan
nhà nước}: dữ liệu thật có nhãn (cắt chủ đề, tiến độ) đều là họp cơ quan nhà nước hoặc Quốc hội; bối cảnh
doanh nghiệp và dự án hiện chỉ được đánh giá trên dữ liệu tổng hợp. \textbf{Dữ liệu đánh giá nhỏ}: 3 cuộc
họp cho cắt chủ đề và 1 cuộc họp cho tiến độ, đều do một người gán nhãn; các khác biệt nhỏ có thể nằm
trong vùng nhiễu. \textbf{Lệch giữa dữ liệu tổng hợp và thật}: bộ cắt chủ đề và ngưỡng của mô hình quyết định
được chọn phần lớn trên dữ liệu tổng hợp; ngưỡng cho phần tiến độ chưa được hiệu chỉnh riêng.
\textbf{Đơn vị là lượt nói}: không tách được lượt nói chứa nhiều chủ đề, và khi bộ tách người nói gộp
nhiều người vào một nhãn thì vai người nói có thể bị hiểu sai. \textbf{Thiên lệch tự đồng ý}: tác tử trả
lời phản hồi chính là tác tử đã trích ra mục đó; dùng mô hình khác họ cho bên kiểm chứng chỉ giảm chứ không
loại bỏ rủi ro này.

\section*{Cân nhắc đạo đức (Ethical Considerations)}

Nội dung họp của cơ quan nhà nước và doanh nghiệp có thể nhạy cảm hoặc là bí mật kinh doanh. Hệ thống nên
chạy trên hạ tầng được phép, ưu tiên mô hình triển khai nội bộ cho các bước xử lý toàn văn, và chỉ ghi âm
khi người tham dự đã được thông báo. \diva{} được thiết kế để hỗ trợ, không thay thế, người soạn và người
duyệt biên bản; biên bản cuối cùng phải do người có trách nhiệm duyệt.

\bibliography{references}

\end{document}
